% =====================================================================
% C7 DRAFT v1.0 (Member C, 2026-10-27) — sec4 + sec5 成稿草稿
% 用法：Part A 复制到 sections/sec4_algorithms.tex，
%       Part B 复制到 sections/sec5_training.tex（替换原 2 行骨架），E 统稿。
% 注意：\cite 键以 references.bib（E 维护）为准；下方已用 TODO 标注待 E 定键。
%       本文件为 C 的草稿暂存，不 \input 进 main.tex，不影响编译。
% =====================================================================

% ===================== Part A → sections/sec4_algorithms.tex =====================
\section{Basic Concepts of Algorithms}
\label{sec:algorithms}

Our main matrix combines four classifiers with four imbalance-handling
strategies. The four classifiers were chosen to span complementary
families of decision-boundary assumptions, and each admits a direct
Weka counterpart used in the comparative discussion of
Section~\ref{sec:training}.

\subsection{Classifiers}
\textbf{Decision Tree (DT).} A CART tree grown with the entropy
impurity criterion and post-pruned by cost-complexity pruning
($\alpha=0.01$); this mirrors Weka's J48 (C4.5), which selects splits
by information gain. Trees yield axis-parallel, rule-like boundaries
and are directly interpretable, but a single tree grown to purity
memorises a 1{,}000-row training set, so light pruning retains
interpretability while limiting overfitting.

\textbf{Random Forest (RF).} An ensemble of 100 de-correlated trees,
each grown on a bootstrap sample of the training fold. Averaging
reduces variance without inflating bias; the resulting boundary is an
ensemble-smoothed version of the tree boundary, and RF tolerates the
class imbalance of both datasets markedly better than any single model
in our plain baselines.

\textbf{Logistic Regression (LR).} A linear decision boundary passed
through a sigmoid; the model is trained by regularised log-loss. LR
produces the best-calibrated probabilities of the four, which matters
for threshold moving (Section~\ref{ssec:strategies}), but the linear
boundary is the most sensitive to how nominal attributes are encoded.

\textbf{Support Vector Classifier (SVC).} A maximum-margin classifier
with an RBF kernel ($C=1.0$), corresponding to Weka's SMO. SVC has no
native probabilities; \texttt{probability=True} wraps an internal
Platt calibration (a 5-fold internal CV fitting a sigmoid to the
decision values), which is what makes AUC and threshold moving
possible at the cost of extra training time.

\subsection{Imbalance-Handling Strategies}
\label{ssec:strategies}
\textbf{plain} keeps the empirical class distribution and the default
0.5 operating point; it is the reference arm for every comparison.
\textbf{SMOTE} over-samples the minority class by interpolating
between minority neighbours and their nearest neighbours; synthesis is
confined to training folds by placing the sampler inside an
\texttt{imblearn} pipeline. \textbf{Cost-sensitive learning} reweights
the training loss: on credit-g we apply the official 5:1 cost matrix
($5\times$FN vs $1\times$FP) via \texttt{class\_weight}; on sick,
which has no official costs, an equal-weight control is used.
\textbf{Threshold moving} keeps training identical to plain and moves
only the operating point: a threshold $\tau$ is selected on inner
5-fold out-of-fold probabilities of the training fold (grid
$0.05{-}0.95$, step 0.05, maximising F1, ties resolved to the smaller
$\tau$), then locked, and the model is refitted on the full training
fold before evaluation. It is the only strategy that never touches
training, hence a clean measure of how much the default 0.5 threshold
leaves on the table.

\subsection{Ensemble Methods: Bagging and Boosting}
\label{ssec:ensemble-concepts}
Beyond the four base classifiers, we conduct an exploratory ensemble
study. \textbf{Bagging} trains multiple instances of a base learner on
bootstrap samples and aggregates their majority votes, reducing
variance; a random forest is bagging combined with random feature
subspaces. \textbf{Boosting} (AdaBoost) instead builds learners
sequentially, each re-weighting the instances the previous ones
misclassified --- an implicit form of cost-sensitive learning on hard
(often minority) examples. Both ensembles use the same pruned decision
tree as the base learner as in Section~\ref{sec:algorithms}, with
$n_{\text{estimators}}=100$ matching the random-forest budget for
comparability.

% ===================== Part B → sections/sec5_training.tex =====================
\section{Model Training}
\label{sec:training}

\subsection{Preprocessing and Pipelines}
Both datasets mix nominal and numeric attributes, and the two model
families require different encodings, so a dual-channel
\texttt{ColumnTransformer} feeds each family what it needs. Trees
(DT, RF) receive nominal attributes through an \texttt{OrdinalEncoder}
--- the integer order is arbitrary but trees only threshold on it, so
no spurious magnitude is introduced --- with numeric attributes passed
through unscaled, as trees are invariant to monotone transforms. The
linear family (LR, SVC) receives \texttt{OneHotEncoder} output (with
\texttt{handle\_unknown='ignore'}) plus \texttt{StandardScaler} on
numeric attributes; feeding ordinal codes into a linear model would
assert false statements such as ``status 2 is twice status 1''.

Every fit-type step --- imputation, encoding, scaling, and SMOTE where
applicable --- lives inside an \texttt{imbalanced-learn} pipeline, so
its parameters are fitted on the training fold only and merely applied
(\texttt{transform}) to the validation fold. On sick, the two fully
missing TBG columns are dropped outside the pipeline (a static,
fold-independent column operation; 100\% missingness is a whole-dataset
fact), while the remaining missing values are imputed with the
training-fold mode inside the pipeline. On credit-g no cleaning is
applied (zero missing values and zero duplicate rows were verified).

\subsection{Locked Configuration}
All hyperparameters were frozen once, before any evaluation, with
\texttt{random\_state=42} on every stochastic component (DT tie-breaks,
RF bootstrap and feature draws, SMOTE neighbour draws, SVC's internal
Platt calibration): DT(\texttt{criterion='entropy'},
$\alpha_{\text{ccp}}=0.01$); RF($n_{\text{estimators}}=100$);
LR($\textit{max\_iter}=1000$, convergence headroom on the 60--80
one-hot dimensions); SVC(RBF, $C=1.0$, \texttt{probability=True}).
Library versions are pinned (scikit-learn 1.9.1, imbalanced-learn
0.14.2, pandas 3.0.6, scipy 1.17.1), so the same versions plus the
locked splits and seed reproduce every reported number bit-for-bit.
% TODO(E): cite scikit-learn (Pedregosa et al., 2011) and
% imbalanced-learn (Lemaitre et al., 2017) once bib keys are fixed.

\subsection{Evaluation Protocol}
Two independent protocols are used. (i) Stratified 10-fold
cross-validation (\texttt{StratifiedKFold(10, shuffle=True,
random\_state=42)}) over the full data, reported as mean$\pm$std over
folds (std with $ddof=0$). (ii) A one-off 70:30 stratified holdout
(700/300 on credit-g; 2640/1132 on sick) evaluated exactly once per
locked configuration. Both splits were generated once, stored as row
index files with a SHA-256 manifest, and are consumed as-is by every
run; no experiment regenerates them. The minority class is the
positive class throughout (bad on credit-g, sick on sick); all metrics
come from a single shared evaluation module (accuracy, positive-class
precision/recall/F1, ROC-AUC, Cohen's $\kappa$, balanced accuracy,
macro-F1, and the credit-g cost $(5\,\mathit{FN}+\mathit{FP})/N$).

\subsection{Dataset-Specific Differences}
Table~\ref{tab:pipeline-diff} summarises how the same locked recipe
specialises to the two datasets.
\begin{table}[htbp]\centering\small
\caption{Two datasets, one locked recipe: differences only.}
\label{tab:pipeline-diff}
\begin{tabular}{lll}
\toprule
Aspect & credit-g & sick \\
\midrule
Rows $\times$ features & $1{,}000\times 20$ & $3{,}772\times 27$ \\
Class ratio (+ : $-$) & 300 : 700 (70:30) & 231 : 3541 ($\approx$94:6) \\
Missing values & none & $\approx$16\%, TBG columns dropped \\
Imputation & --- & mode, in-pipeline \\
Cost-sensitive arm & official 5:1 \texttt{class\_weight} & equal-weight control \\
Holdout split & 700 / 300 & 2640 / 1132 \\
\bottomrule
\end{tabular}
\end{table}

\paragraph{Ensemble study setup}
The ensemble configurations (Bagging/AdaBoost $\times$ plain/SMOTE,
both protocols, both datasets) reuse the locked cross-validation splits
and the unified \texttt{evaluate.py} pipeline of the main matrix, but
are stored in a separate result namespace (\texttt{results/ensemble/})
and are \emph{not} merged into the 32-configuration matrix, so that the
frozen statistical comparison remains untouched. Runtimes are recorded
per fold by the same runner.
